% Figaro Paper D: FIG — A Schelling-Point Token for the Figaro Coordination Ecosystem
% Cryptoeconomic design. Snapshot of the implementation as of 2026-04-23.
% The kernel itself is token-agnostic and is treated in figaro3a.tex / figaro3c.tex.

\documentclass[11pt,a4paper]{article}

\usepackage[utf8]{inputenc}
\usepackage[T1]{fontenc}
\usepackage{amsmath,amssymb,amsthm}
\usepackage{mathtools}
\usepackage{booktabs}
\usepackage{graphicx}
\usepackage{hyperref}
\usepackage{cleveref}
\usepackage{enumitem}
\usepackage{xcolor}
\usepackage[margin=1in]{geometry}
\usepackage{natbib}
\bibliographystyle{plainnat}

\newtheorem{theorem}{Theorem}[section]
\newtheorem{proposition}[theorem]{Proposition}
\theoremstyle{definition}
\newtheorem{definition}[theorem]{Definition}
\theoremstyle{remark}
\newtheorem{remark}[theorem]{Remark}

\newcommand{\figaro}{\textsc{Figaro}}
\newcommand{\fig}{\textsc{Fig}}
\newcommand{\pay}{P}

\title{%
  \fig{}: A Schelling-Point Token\\
  for the Figaro Coordination Ecosystem%
}

\author{
  Alessandro Daliana\thanks{Corresponding author. \texttt{adaliana@figaro.org}}
}

\date{April 2026 (snapshot: 2026-04-23)}

\begin{document}

\maketitle

% ════════════════════════════════════════════════════════════════════
\begin{abstract}
We describe \fig, the coordination token that ships alongside the
\figaro{} settlement primitive. \fig{} is designed as a
\emph{Schelling-point token} in the sense of
\citet{schelling1960strategy}: a focal point around which
participants in the \figaro{} ecosystem can coordinate without
central authority or explicit agreement. Its value as a
coordination artifact derives from its focality, not from any
on-chain utility function, yield stream, governance claim, or fee
path.

This framing inverts the usual design question for a new token.
Rather than asking ``what does FIG \emph{do}?'' we ask ``what
would \fig{} have to \emph{not} do in order to remain a credible
focal point?'' The answer produces a catalogue of design
exclusions, each of which removes a class of holding motivation
that would dilute the purity of \fig{}'s signaling value. \fig{}
is not governance, not yield-bearing, not a fee path, not bond
collateral for the kernel, not settlement-emitted, and not
founder-vested. Each negative preserves focality by removing a
reason to hold \fig{} that is not ``I am signaling alignment with
the Figaro ecosystem.''

The positive design content is modest and targeted: a fixed
one-billion-unit cap enforced by two complementary supply-integrity
mechanisms, a one-way deployer-renouncement latch that freezes the
allocation plan at deployment, and a single staged-merkle airdrop
that distributes 60\% of supply to pre-committed cohorts at years
2, 5, and 9. We lay this design out in detail, reject an earlier
settlement-anchored ``Euler oscillation'' design for the
focality-destroying reasons it introduced, and close with an
honest analysis of how \fig{} can lose its focal-point status and
what the design does and does not defend against.

The paper sits within a lineage. \citet{nakamoto2008bitcoin}
demonstrated that money could be a protocol rather than an
institution, with \textsc{Btc} as the focal Schelling point for a
peer-to-peer cash system. \citet{buterin2014ethereum} demonstrated
that contracts could be composable, with \textsc{Eth} as the focal
point for programmable-settlement. The present work extends the
lineage one step: \emph{institutions} themselves become composable
through the \figaro{}
primitive (mechanism and implementation papers), and
\fig{} is the focal point for participants who coordinate around
the composition.
\end{abstract}

\medskip
\noindent\textbf{Keywords:}
cryptoeconomics, Schelling point, focal-point coordination, token
design, coordination token, supply integrity, decoupled emission

% ════════════════════════════════════════════════════════════════════
\section{Introduction}
\label{sec:intro}

\fig{} is the native token that ships alongside the \figaro{}
settlement primitive developed
in the mechanism and implementation papers. This paper
describes its design and defends a specific framing of what it
is. The framing is the substantive contribution; the
supply-integrity and distribution mechanics follow from it.

The framing: \emph{\fig{} is a Schelling-point token for the
\figaro{} ecosystem.} Its value to participants is not a claim on
protocol revenue (there is no protocol fee), not a governance
vote (the kernel has no governance), not a yield stream (no
activity-coupled emission), not a utility token purchasable for
on-chain services (there are no protocol-layer services that
require it), and not bond collateral (the kernel is
token-agnostic). Its value is its focality: in the sense of
\citet{schelling1960strategy}, \fig{} is the focal point that
\figaro{}-aligned participants converge on when they want to
signal alignment without coordinating through a central authority.

\paragraph{A lineage note.}
The design pattern has precedent. \citet{nakamoto2008bitcoin}'s
Bitcoin is, in retrospect, the clearest example of a
Schelling-point token at the money layer: its value as a store of
value derives from holders converging on it as \emph{the} answer
to the question ``what digital bearer asset are we
coordinating around?'' The proof-of-work mechanism secures
Bitcoin's ledger, but it does not manufacture Bitcoin's
Schelling-point status --- that emerged from participant
convergence under conditions in which Bitcoin's design made
convergence viable (fixed supply, open participation, no issuer).
\citet{buterin2014ethereum}'s Ethereum moved the logic up one
layer: \textsc{Eth} became the focal point for programmable
settlement, with composable contracts as the substrate. The
present work moves it up again: with \figaro{}, institutions
themselves are composable, and \fig{} is the focal point for
aligned participation in that composition. Each layer's native
token has the same structural role --- a focal coordination
anchor that does not require a central issuer, does not require a
governance apparatus, and cannot be replicated by a
general-purpose token, because the signal it provides is
specifically about alignment with that particular ecosystem.
Each layer's token is also underpinned by a different primary
economic quantity --- Bitcoin by the cost of producing the asset,
Ethereum by the demand for executing transactions, \fig{} by the
trade conducted through the substrate (\Cref{sec:ecosystem-need}).
Of the three, \fig{}'s underpinning sits closest to where actual
economic value is created, because trade is where it is created.

\paragraph{Paper organization.}
\Cref{sec:design-category} introduces Schelling-point tokens as a
design category and locates \fig{} within it.
\Cref{sec:ecosystem-need} argues why the \figaro{} ecosystem is
well-suited to a focal-point native token rather than a
utility-coupled one.
\Cref{sec:design-rules} articulates design rules that support
focal-point emergence and resist focality-dilution.
\Cref{sec:supply} develops supply integrity.
\Cref{sec:allocation} develops allocation.
\Cref{sec:vesting} develops the staged-merkle airdrop.
\Cref{sec:negatives} catalogues the design exclusions and ties
each to the focality property it preserves.
\Cref{sec:rejected} records the rejection of an earlier
settlement-anchored ``Euler oscillation'' design and deepens the
rejection rationale in light of the Schelling-point framing.
\Cref{sec:bootstrap} treats bootstrap and attack surface.
\Cref{sec:silhouette} retains the scope note on
substrate-neutrality and distributional silhouette.
\Cref{sec:verification} reports the verification surface.
\Cref{sec:conclusion} concludes.

% ════════════════════════════════════════════════════════════════════
\section{Schelling-Point Tokens as a Design Category}
\label{sec:design-category}

\citet{schelling1960strategy} identified \emph{focal points}
(what later came to be called Schelling points) as solutions that
parties tend to choose in the absence of explicit communication,
because some feature of the solution makes it stand out as the
natural coordination target. Schelling's canonical example: if
two parties are told to meet in New York City without being able
to communicate, a disproportionate fraction converge on Grand
Central at noon. Neither the location nor the time is more
rational than any other; what makes them focal is a shared
cultural salience that each party correctly expects the other to
recognize.

\paragraph{The focal-point literature, briefly.}
The conceptual machinery is older and more developed than the
crypto-token literature usually engages. \citet{schelling1960strategy}
introduced focal points as a coordination phenomenon and observed
the empirical fact that humans converge on salient solutions
without communication. \citet{sugden1995theory} formalized the
phenomenon as \emph{salience theory}: parties pick the option that
stands out by some shared frame of reference, and the salience is
itself constructed through repeated coordination on similar
problems. \citet{mehta_starmer_sugden_1994} produced the canonical
experimental confirmation that focal-point convergence operates
robustly across populations on a wide range of pure-coordination
problems. \citet{bacharach2006beyond}'s \emph{team reasoning}
formalized the further step that focal points operate not because
each party reasons about the other's reasoning but because both
parties shift to a we-frame in which the salient solution is the
team-rational choice. The combined apparatus gives a precise
account of what makes a token focal: it is not the token's
properties in isolation but the token's properties as picked out
by a shared frame within which holding the token is recognizable
as a coordination act.

The cryptoeconomic literature has applied the apparatus
repeatedly without always citing it.
\citet{markey-towler2018schelling} examines tokens as Schelling
points in a behavioral-economics register; \citet{daian_etal_2017}
and the broader Robust Incentives Group treat focal-point
properties as design objectives in mechanism analysis;
\citet{buterin2020credible} on \emph{credible neutrality}
articulates the design discipline for tokens whose focality is to
be sustained over time. The classical economic-theory antecedents
worth naming are \citet{hayek1945knowledge} on the
distributed-coordination role of prices (a non-token analogue of
the focal-point function) and \citet{selgin1988theory} on the
emergence of focal monetary instruments without central authority.

\paragraph{Crypto and Schelling points.}
The structure has been rediscovered repeatedly in crypto design.
\citet{nakamoto2008bitcoin} did not use the word ``Schelling,''
but Bitcoin's emergence as a store of value followed the logic
precisely: a fixed-supply, non-upgradable, issuer-less digital
bearer asset with a peer-to-peer monetary intent. No mechanism
guarantees Bitcoin's SoV status; it emerged because the design
made focal convergence viable and participants converged.
\citet{buterin2014ethereum} extended the logic up a layer:
\textsc{Eth} became the focal settlement unit for programmable
contracts, for the same reason --- the design made convergence
viable (gas-required, ubiquitous across contracts, issuer-less at
the base layer). Similar focal patterns have emerged for
subsidiary positions within crypto: Uniswap as the focal AMM,
Aave as the focal lending protocol, Lido as the focal liquid-staking
wrapper. In each case the token (\textsc{Uni},
\textsc{Aave}, \textsc{Ldo}) has some utility, but a substantial
fraction of its holding behavior is Schelling convergence on
``the token that represents alignment with this particular
primitive.''

\paragraph{Schelling tokens vs.\ utility tokens.}
A pure utility token is valued for the services a holder can
purchase with it (gas, access, data). A pure governance token is
valued for the votes it conveys. A pure yield-bearing token is
valued for the cash flows it represents. A Schelling-point token
is valued for something different: its \emph{focality}, which is
a coordination property, not a cash-flow property or a rights
property. A participant holds a Schelling-point token because
holding it makes them legible to other aligned participants and
because the shared holding constitutes the coordination itself.

Pure Schelling-point tokens are rare because most token designs
add utility or governance features to create additional demand
drivers. Each such addition, however, introduces holders whose
reason for holding is \emph{not} focal alignment --- governance
speculators, yield farmers, utility arbitrageurs --- and these
holders' presence dilutes the focality signal. Bitcoin's focality
is durable partly because Bitcoin has \emph{remained} a pure
Schelling point by refusing to add governance or yield or utility
features beyond the base-layer payment rail. Each refused feature
protects the signal's purity.

\paragraph{\fig{}'s design category.}
\fig{} is designed to sit in the pure-Schelling-point category.
The design is committed to refusing utility couplings, governance
rights, yield streams, and fee paths. The rest of this paper
develops what that refusal looks like in practice.

% ════════════════════════════════════════════════════════════════════
\section{Why the \figaro{} Ecosystem Is Well-Suited to a Focal-Point Token}
\label{sec:ecosystem-need}

\paragraph{What underpins \fig{}'s value, said honestly.}
A Schelling-point token requires focal convergence, and focal
convergence is itself a value-formation mechanism in the
strict sense: a holder values the token because other aligned
holders value it for the same reason, and the shared
recognition is what gives the token its content. We need to be
careful here, because earlier framings of this paper (and the
adjacent crypto literature broadly) have invited a misreading
that conflates focal convergence with mechanism-coupled token
demand. The two are different and should be kept apart.

The standard lineage --- Bitcoin underpinned by the cost of
proof-of-work, Ethereum underpinned by gas demand, FIG
underpinned by trade --- is rhetorically tidy and mechanically
misleading. Bitcoin and Ethereum each have a
\emph{mechanism-coupled} demand source: Bitcoin's PoW
issuance consumes real factors and puts a marginal-cost floor
under issuance; Ethereum's gas requirement makes \textsc{Eth}
necessary for any transaction on the Ethereum L1, with
EIP-1559 burn reinforcing the floor. \fig{} has neither.
The \figaro{} kernel uses any ERC-20 as its bond and payment
currency; a participant who settles every bonded commitment
in \texttt{USDC} pays zero \fig{}, contributes zero
\fig{}-denominated demand, and the kernel is indifferent to
their choice. Saying \fig{} is ``underpinned by trade'' in
the way Bitcoin is underpinned by proof-of-work elides this
asymmetry and overstates the case.

\fig{} is a \emph{pure focal-point token} in the sense
\Cref{sec:design-category} defined. Its value derives from
the focal-convergence mechanism alone --- aligned holders
recognize \fig{} as the coordination anchor for the
\figaro{}-ecosystem question, and the coordination is itself
the token's content. There is no settlement coupling, no gas
coupling, no fee path that routes ecosystem activity through
\fig{}-denominated demand. We do not introduce one, and the
design rules of \Cref{sec:design-rules} below explicitly
forbid us from doing so.

This is a more honest version of the standard lineage claim,
and we think it is also the correct version. \fig{} occupies
a different category from Bitcoin and Ethereum, not a
homologous position in the same category. Bitcoin and
Ethereum are tokens whose mechanism-couplings happen to
support focal convergence on top; \fig{} is a token whose
focal convergence is the mechanism, with no coupling beneath
it. The question is whether such a category can sustain
value, and the answer is uncertain in the way any pure
coordination phenomenon is uncertain: focality can be
sustained through participant convergence, can be diluted by
holders whose reason to hold is not focal alignment, and can
be lost altogether if the convergence breaks. The design
discipline below is the work of keeping the convergence
intact.

A reader sympathetic to the cost-of-production or
transaction-demand framings should understand the present
paper as conceding their mechanical correctness and arguing
that they do not apply to \fig{}, rather than claiming
\fig{} sits in the same category. The trade-volume of the
\figaro{} ecosystem is a useful public signal of the
ecosystem's health, and a holder of \fig{} may correctly
expect that ecosystem health correlates with the maintenance
of focal alignment, but the correlation is reflexive (held
together by participant belief) rather than mechanical (held
together by required-token-flow). Reflexive value is real
and durable when the focality is durable; it is not the same
thing as cost-floor or fee-floor value.

\paragraph{Why a general-purpose token cannot carry this signal.}
The \figaro{} kernel is token-agnostic. Any ERC-20 can serve as
the bond and payment currency for a bonded
commitment (mechanism paper). The kernel does not need
a native token to function, and \fig{} does not interact with
the kernel in any equilibrium-relevant way. Why, then, have a
native token at all?

The answer is that multi-party coordination in an ecosystem
without central authority requires a coordination signal that
does not depend on any central authority, and no general-purpose
token provides that signal. Participants bond in \texttt{USDC}
when they want legal-system alignment; in a DAO governance token
when they want community membership in that DAO; in \texttt{ETH}
when they want to denominate in the base settlement unit; in a
commodity-backed token when they want value anchoring. None of
these signals \emph{Figaro-ecosystem alignment specifically}. A
general-purpose token carries many signals; what distinguishes
them cannot be inferred from the token.

A focal-point token that signals Figaro alignment \emph{only},
and whose design forbids any other reading of the holder's
motivation, provides a coordination instrument the ecosystem does
not otherwise have. This instrument is useful in two distinct
settings:

\paragraph{Finding aligned counterparties.}
A participant designing a new \figaro{} assembly (a bonded
workflow composed from protocol components) wants to find
counterparties who understand the protocol, operate in its idiom,
and will engage constructively with its equilibrium discipline.
Searching for such counterparties by settlement history is
possible (§5 of the institutional-economics paper) but expensive --- one must
read the chain. Searching by \fig{} holding is cheaper: a
participant who holds \fig{} has declared alignment through the
specific act of acquiring a token whose only plausible reason to
acquire is alignment.

\paragraph{Signaling alignment in contested contexts.}
Participants who want to declare \figaro{}-aligned participation
in adjacent debates --- regulatory, academic, industry --- benefit
from a token whose holding is unambiguous. A founder who accepts
\fig{} in lieu of other compensation signals more than a founder
who accepts \texttt{USDC}. Not every participant needs this
signal; those who do benefit from the existence of a focal
coordination anchor that has no financial pretext.

These use cases are modest. We do not claim \fig{} is
indispensable or that the ecosystem fails without it. We claim
only that a focal-point token has a coordination function
unavailable to general-purpose tokens, and that a cohort of
participants will find the function useful.

% ════════════════════════════════════════════════════════════════════
\section{Design Rules for a Focal-Point Token}
\label{sec:design-rules}

Supporting focal-point emergence is not the same as manufacturing
it. A designer cannot declare that a token is a Schelling point;
focality emerges from participant convergence or it does not.
What the designer can do is remove features that would make
convergence implausible or fragile, and add features that make
convergence stable once it starts.

We articulate four design rules.

\paragraph{Rule 1: Pure signaling.}
The token must carry \emph{one} unambiguous reason to hold it.
Holders should be alignment-signalers, not yield-chasers,
governance-speculators, utility-arbitrageurs, or
fee-path-optimizers. Features that would introduce non-alignment
holders are forbidden. This is the most restrictive rule and
generates most of \fig{}'s catalogue of exclusions
(\Cref{sec:negatives}).

\paragraph{Rule 2: Predictable supply.}
Focality is fragile under surprise inflation. A holder who
believes supply is fixed may exit on learning it is not.
\fig{}'s supply schedule is committed at deployment and cannot
be modified: 1B hard cap, staged airdrop with three immutable
merkle roots and immutable calendar unlocks, one-way
deployer-renouncement at the end of the deploy transaction. There
are no surprise-inflation paths available to any actor after
deployment.

\paragraph{Rule 3: Clean distribution.}
Focality is fragile under highly concentrated ownership. If a
single holder can dump enough supply to destabilize the price,
Schelling convergence is less credible. \fig{}'s allocation is
60\% to a pre-committed community cohort (the staged airdrop),
30\% to a treasury with an ecosystem mandate (the DAO), and 10\%
to founders. The concentration is higher than ideal for
focality (the DAO + founder share is 40\%), which we address
honestly in \Cref{sec:bootstrap} and in the
distributional-silhouette concession in
\Cref{sec:silhouette}.

\paragraph{Rule 4: Absence of protocol-activity coupling.}
Focality requires that participants hold for alignment reasons.
If the token's supply is a function of protocol activity,
holders acquire the token by doing things that are not alignment
(gaming the emission curve, wash-trading for rewards), and the
signal is muddied. Activity-coupled emission is ruled out. We
develop the argument at length in \Cref{sec:rejected}, where we
describe an earlier \fig{} design that used settlement-anchored
emission and explain why we rejected it on focality grounds.

% ════════════════════════════════════════════════════════════════════
\section{Supply Integrity}
\label{sec:supply}

\fig{} has a fixed maximum supply of $10^9$ units (one billion,
denominated in 18-decimal $\mathit{ether}$ units in the
contract). The supply ceiling is enforced by two distinct
mechanisms operating at different times.

\paragraph{Per-mint enforcement (run-time).}
The \texttt{mint} function reverts with \texttt{SupplyCapExceeded}
if $\texttt{totalSupply()} + \texttt{amount} > \texttt{MAX\_SUPPLY}$,
checked atomically inside a non-reentrant call.

\paragraph{Pre-registration enforcement (deploy-time).}
The contract maintains an aggregate \texttt{totalRegisteredCap}:
the sum of all registered minter caps. \texttt{registerMinter}
reverts if
$\texttt{totalRegisteredCap} + \texttt{cap} > \texttt{MAX\_SUPPLY}$.
This means the entire allocation plan is committed to the
$10^9$ ceiling \emph{before any token is minted}.

The pair (per-mint + pre-registration) is stronger than either
alone in the focality frame: per-mint enforcement catches cap
overruns that survive bugs; pre-registration enforcement puts the
allocation plan on chain where a focality-assessing participant
can read it directly rather than relying on off-chain narrative.

\paragraph{The renouncement latch.}
\fig{} has a one-way Boolean \texttt{deployerMintRenounced},
settable exactly once by the original deployer via
\texttt{renounceDeployerMint}. After renouncement,
\texttt{registerMinter} reverts and any deployer call to
\texttt{mint} reverts. The deploy script
(\texttt{script/DeployMainnet.s.sol}) registers exactly two
minters --- the deployer wallet (cap $4 \times 10^8$, used for
the genesis mint) and the airdrop contract (cap $6 \times
10^8$) --- and then calls \texttt{renounceDeployerMint} as the
final action of the deploy transaction. After the transaction
completes:

\begin{itemize}
  \item $\texttt{totalRegisteredCap} = 10^9$ exactly;
  \item the deployer has minted exactly $4 \times 10^8$ to the
    designated founder and DAO wallets;
  \item only the airdrop contract retains residual mint
    authority, capped at $6 \times 10^8$;
  \item no path exists to register additional minters.
\end{itemize}

The supply schedule is therefore frozen by a single transaction,
and the contract has no further ``governance moments'' available
to anyone --- not to the deployer, not to the DAO, not to any
future upgrade proposal. In focality terms: a participant
assessing \fig{} at time $t > t_{\text{deploy}}$ can verify on
chain that no further minting will occur beyond the airdrop
schedule. Focal convergence is not destabilized by uncertainty
about supply.

% ════════════════════════════════════════════════════════════════════
\section{Allocation: 10 / 30 / 60}
\label{sec:allocation}

The total supply is split:

\begin{center}
\begin{tabular}{lrl}
\toprule
\textbf{Bucket} & \textbf{Share} & \textbf{Realization} \\
\midrule
Founder      & 10\% (100M)  & Genesis mint, no vesting, no unlock \\
DAO          & 30\% (300M)  & Genesis mint, no vesting, no unlock \\
Airdrop      & 60\% (600M)  & Three-stage merkle, yr 2 / 5 / 9 \\
\bottomrule
\end{tabular}
\end{center}

\paragraph{Founder allocation (10\%).}
A modest founder share, by contemporary norms. The non-vesting
choice is deliberate: a vesting cliff is a tacit promise of
future voting alignment or schedule-driven holding, both of which
would be reasons to hold \fig{} other than alignment signaling.
Releasing the founder allocation at genesis exposes the founder
to whatever market consequences attach to early sale; the
structural observation is that an unvested founder forgoes an
asymmetric option (sell into early demand if it materializes;
ride out the vesting if it does not) that a vested founder
retains, and we have chosen to forgo the option. The prior
question --- whether founder allocation at this magnitude is the
right allocation at all --- is a political-economy question we
do not attempt to answer in a cryptoeconomic paper.

\paragraph{The bootstrap-dump scenario.}
A fair objection to the no-vesting choice is the bootstrap-dump
scenario: between genesis (year 0) and the year-2 first airdrop
stage, the founder cohort holds a 10\% liquid position and the
focal-point convergence has not yet had time to establish the
Schelling stickiness that would make later dump events
self-correcting. A founder dump in this window could foreclose
focality emergence by signaling to potential aligned holders
that the alignment is not credibly committed. The §8 attack-surface
analysis treats general dump attacks but does not specifically
address the bootstrap window where the structural defenses are
weakest.

The honest answer is that the design has accepted this risk
rather than mitigated it through vesting. Two reasons. First,
vesting would import the same focality-dilution problem the
non-vesting choice was designed to avoid: a founder who must
vest will be a holder for schedule-driven reasons that are not
alignment, and the schedule itself becomes a Schelling object
that competes with the underlying alignment signal. Second, a
public lock-up commitment outside a vesting contract --- e.g.,
self-custody with stated non-sale intent for a specified
window --- is available to a founder who wishes to mitigate the
bootstrap risk through credible signal rather than through
mechanism. We have not committed to a specific lock-up window
in the present paper because the signal's content depends on
the founder's standing relationship with the participant
cohort rather than on a contractual instrument; the
appropriate commitment is a matter for the founders to make
publicly and is not a property of the protocol itself.

A reader skeptical of this trade-off should weight the
bootstrap-dump risk against the focality-dilution risk and
make their own judgment about which is structurally worse. We
think the focality-dilution risk is the larger one because it
is irreversible in a way the bootstrap-dump risk is not (a
dump can be survived if alignment is durable; a focality
diluted by vested-holder noise cannot be recovered without
restarting), but we do not pretend the trade-off is settled.

\paragraph{DAO allocation (30\%).}
A treasury share, available immediately at genesis. The DAO is
not a kernel-governance body --- there is no kernel governance to
be done --- but a coordination treasury for ecosystem activity
(grants, public goods, reference implementations). Releasing the
DAO allocation at genesis means the treasury is operational from
day one. A fair objection: doesn't a liquid 30\% treasury with
an ecosystem-grants mandate create exactly the reflexivity that
\Cref{sec:rejected} rejects in the Euler design, and doesn't it
dilute \fig{}'s focal-point purity (Rule 1)? We concede the
objection in its general form. The defense is narrower than the
objection implies: the coupling the Euler design introduces is
\emph{mechanism-level} (a deterministic on-chain function from
settlement count to mint), while the coupling a treasury
introduces is \emph{policy-level} (a human-discretionary
function from treasury governance to disbursement).
Mechanism-level coupling is what the kernel's no-escape-hatches
design principle rules out; policy-level coupling is attenuated
by two layers of friction (off-chain discretion, revocable
grants). The distinction is principled but not absolute: a
sufficiently mechanical grant policy would reintroduce the
coupling, which is why the DAO's operating procedures are
themselves a design artifact that warrants its own treatment
rather than a silent assumption.

\paragraph{Airdrop allocation (60\%).}
The majority share, allocated over a calendar window with
declining per-stage size. The structure is designed to (a) reach
non-genesis participants, (b) reward sustained participation
(longer-window cohorts at year 5 and year 9), and (c) avoid the
``dump-on-launch'' dynamic of single-event airdrops by spreading
the supply emission over a 9-year horizon.

The 60/30/10 internal split of the airdrop (300M at yr 2, 200M
at yr 5, 100M at yr 9) is front-loaded for the bootstrap
reason: the network needs participants more in its early years
than in its mature years. The declining schedule encodes that
asymmetry into the supply.

\paragraph{A note on the specific percentages.}
The 10/30/60 split and the internal 50/33/17 staging of the
airdrop reflect a design judgment rather than a first-principles
derivation. Any front-loaded schedule (45/35/20, 40/40/20,
50/30/20) would satisfy the bootstrap argument above, and the
specific weights are a calibration we do not attempt to justify
beyond ``plausible under the stated principles.''

% ════════════════════════════════════════════════════════════════════
\section{Cohort-Based Vesting via Staged Merkle}
\label{sec:vesting}

The 600M airdrop is realized by a single contract,
\texttt{RpgfMinter.sol}, with three immutable stages
($\texttt{STAGE\_COUNT} = 3$). Each stage is parameterized by an
immutable \texttt{unlockTime}; the per-stage merkle \texttt{root}
is submitted at tranche time by a sequencer-role address, gated
by verification of an SP1 zero-knowledge proof attesting that the
root correctly aggregates the schema-author substrate-broadening
formula. Per-stage root submission is one-shot --- the root
cannot be replaced once committed. There is no admin-controlled
stage addition, schedule extension, or cohort modification.

(An earlier candidate design baked the three merkle roots into
the constructor alongside the unlock times under a contract named
\texttt{StagedMerkleAirdrop.sol}; that shape was retired in
2026-05 when the schema-author retroactive-funding aggregation
matured to the point where pre-deploy root computation became
inadequate.)

\paragraph{Single contract, three stages.}
An earlier candidate design used three separate airdrop
contracts. The single-contract design is preferred for two
reasons. First, surface area: a single contract is a single
audit target, a single set of formal-verification rules, and a
single storage account. Second, atomic deploy guarantees: the
\texttt{registerMinter} call binds the airdrop contract's $6
\times 10^8$ cap before the renouncement, and a single contract
makes the cap binding atomic.

\paragraph{Leaf format.}
The merkle leaf is
$\texttt{keccak256(abi.encodePacked(account, amount))}$, the
canonical OpenZeppelin format. Choosing a non-canonical leaf
format would have required custom tooling per off-chain proof
generator and a corresponding audit cost; we accept the canonical
choice.

\paragraph{One-shot per (stage, address).}
The \texttt{claimed[stageIndex][msg.sender]} double-mapping is
set true on successful claim and checked on every entry. A given
account can claim at most once per stage. No top-up, no
re-claim, no rebate.

\paragraph{Calendar-based unlock.}
The unlock timestamps are calendar-based (Unix epoch seconds),
not activity-based. Tying unlocks to settlement activity would
re-introduce a coupling between protocol use and token
availability --- the category Rule 4 forbids. Calendar timing
has the additional virtue of being predictable to all parties at
deploy time, which supports focal-point clarity.

% ════════════════════════════════════════════════════════════════════
\section{Design Exclusions as Focality Preservers}
\label{sec:negatives}

The cryptoeconomic content of \fig{} is largely defined by what
the design excludes. We catalogue the exclusions and tie each
explicitly to the focality property it preserves. This restates
the exclusions documented in earlier drafts of the paper, but
with the Schelling-point framing the reasons become unified: each
exclusion removes a class of holder whose reason for holding
\fig{} would be something other than alignment signaling.

\paragraph{\fig{} is not a governance token.}
The kernel has no governance surface to vote on. Adding a
governance role to \fig{} would introduce governance
speculators --- participants who hold \fig{} because they want
to vote on proposals, not because they want to signal
alignment. This would dilute the signal by populating the
holder set with non-alignment holders. Refusing governance
preserves the signal's purity (Rule 1).

\paragraph{\fig{} is not yield-bearing.}
The kernel takes no protocol fee; there is no fee pool to
distribute. A stake-and-earn mechanism would introduce yield
farmers --- participants who hold \fig{} for expected cash flow,
not for alignment. Refusing yield preserves the signal's
purity (Rule 1) and also preserves the activity-decoupling
(Rule 4), since yield mechanisms are typically indirectly coupled
to protocol activity.

\paragraph{\fig{} is not a fee path.}
Even if a future runtime layer above the kernel introduces a
fee mechanism, that fee is not denominated in \fig. A
fee-denomination role would introduce utility arbitrageurs ---
participants who acquire \fig{} to pay fees, not to signal
alignment. The signal would be muddied by transactional
holders.

\paragraph{\fig{} is not bond collateral for the kernel.}
The kernel's bonds are denominated in whatever ERC-20 the
parties choose. Making \fig{} the preferred bond denomination
would introduce collateral demand that distorts price discovery
and couples \fig{} to kernel activity through the back door.
Refusing this role keeps the kernel token-agnostic (per Paper
A's design) and preserves \fig{}'s focal-signal purity.

\paragraph{\fig{} is not founder-vested.}
The 10\% founder allocation and the 30\% DAO allocation are
unvested. There is no cliff, no schedule, no time-locked vesting
contract. Vesting would introduce schedule-driven holders ---
participants who hold \fig{} because they will receive more in
the future, not because they are signaling current alignment.
The vested-founder case is the clearest instance: a founder
vesting into liquidity over four years holds \fig{} during that
period for mixed reasons, only one of which is alignment.

\paragraph{\fig{} is not settlement-emitted.}
There is no path from kernel activity (commit, resolve, attest)
to \fig{} mint. \texttt{FigaroBatchVerifier} --- the layer-2
batch verification contract --- is explicitly not a registered
minter and will never be. The deploy script does not register
it; the renouncement latch ensures it cannot be added later.
This refusal is the most rigorous application of Rule 4 and is
worth its own section (\Cref{sec:rejected}).

\paragraph{\fig{} is not a coordinator-pattern extension.}
A \emph{coordinator pattern} for protocol composition specifies
four sufficient conditions under which a mechanism composed with
the kernel preserves the bonding equilibrium. \fig{} satisfies these
conditions trivially: it has no kernel interaction at all.
\fig{} and the kernel are siblings, not a composition.

% ════════════════════════════════════════════════════════════════════
\section{Rejected Design: Euler-Oscillation Emission}
\label{sec:rejected}

An earlier iteration (preprint draft, February 2026) included
a settlement-anchored emission mechanism we called \emph{Euler
oscillation}: a settlement-count-decayed cosine modulation on
top of a halving envelope, with seller-only emission at
resolution time. Three properties motivated it: persistent
bootstrap subsidy, partial wash-trade resistance, and seasonal
coordination through a $1.5\times$ peak rate.

We rejected it. The rejection deepens under the Schelling-point
frame. In the prior draft we listed three reasons: reflexivity,
partial-but-not-eliminative wash-trade defense, and
oscillation-pattern speculation. The Schelling-point frame
reveals a more fundamental reason: activity-coupled emission
\emph{destroys focal-point status} by introducing holders whose
reason for acquiring \fig{} is not alignment.

\paragraph{The focality argument, formally.}
Under Euler emission at rate $r(n)$ with settlement-count decay
and seller-only payout, a participant who settles a process
through the kernel acquires \fig{}. The acquisition is
triggered by the settlement event, not by a decision to signal
alignment. After acquisition, the holder has \fig{} in their
wallet --- but the token-in-wallet does not distinguish
``holder of \fig{} by alignment choice'' from ``holder of
\fig{} because an Euler function happened to pay me.'' In the
Schelling-point frame, these are two distinct populations whose
signals interfere. A third party looking at the distribution of
\fig{} holders cannot read focal alignment from the
distribution because the distribution mixes two holder
populations.

This is distinct from the reflexivity argument (which points at
a decision-coupling between settlement choice and emission
expectation) and the wash-trade argument (which points at a
gameable payoff structure). It is a more fundamental point about
what a token \emph{means} under each design. Activity-coupled
emission makes the token mean two things at once; pure airdrop
distribution makes it mean one thing.

\paragraph{The wash-trade inequality, for completeness.}
Let the prevailing \fig{}/currency exchange rate be $\phi$, let
the seller-only emission rate at settlement count $n$ be $r(n)$,
let gas per cycle be $g$, and let the process duration be $T$
with opportunity cost of capital $\rho$ per unit time. A
wash-trader bonds both sides: buyer posts $2\pay$, seller posts
$2\pay$, total locked $4\pay$ for $T$, seller receives $r(n)$
\fig{} at resolution. The cycle is profitable iff
\begin{equation}
  \phi \cdot r(n) \;>\; \rho \cdot 4\pay \cdot T \;+\; g.
  \label{eq:wash-trade}
\end{equation}
Seller-only emission doubles the capital requirement from
$2\pay$ to $4\pay$ but does not drive \eqref{eq:wash-trade} to
infeasibility. The parameter region where wash-trading is
infeasible is non-empty but not universal. This was a secondary
reason to reject Euler; the primary reason was the
focality-destroying mixing of holder motivations above.

\paragraph{What we kept.}
The 600M airdrop bucket, the staged structure (now realized via
calendar unlocks rather than settlement counts), and the design
intent of distributing the bootstrap subsidy over time. The
underlying intuition --- that distribution should be staged
rather than instantaneous --- survived; the Euler mathematics
did not.

% ════════════════════════════════════════════════════════════════════
\section{Bootstrap and Attack Surface}
\label{sec:bootstrap}

A Schelling-point token cannot be declared; it must be
\emph{recognized}. This section is the most uncertain part of
the design, because it concerns processes we cannot directly
control.

\paragraph{Bootstrap: how \fig{} achieves focal status.}
Focal-point emergence for a newly-launched token depends on
several converging factors, none of which the designer can
guarantee.

\begin{itemize}
  \item \emph{Initial cohort.} The year-2 airdrop distributes
    300M to a pre-committed cohort determined by a merkle root
    set at deploy. The cohort selection process (which is
    policy, not mechanism) determines who receives \fig{} in
    the first distribution and therefore who has the first
    opportunity to recognize \fig{} as a focal point.
  \item \emph{Design legibility.} The design rules of
    \Cref{sec:design-rules}, and the exclusions of
    \Cref{sec:negatives}, must be legible to prospective
    holders as consistent with focal-point coordination rather
    than as random design preferences. This paper is part of
    that legibility work.
  \item \emph{Absence of competing focal candidates.} If an
    alternative token launched with better focal properties
    (cleaner distribution, more restrictive exclusions) at a
    moment when \fig{} had not yet stabilized, participants
    might converge on it instead. The design attempts to make
    \fig{} hard to beat on focal grounds (Rule 1 purity,
    Rule 2 supply predictability) but cannot preempt all
    competitors.
\end{itemize}

\paragraph{Focal-shift attack.}
The relevant attack against a Schelling-point token is the
\emph{focal-shift attack}: an adversary, through concerted
action, moves focal convergence to a different token. The
canonical example is the decade-long debate over whether
Bitcoin-Cash or Bitcoin-SV could displace Bitcoin as the SoV
focal point; neither succeeded, but both made the attempt at
significant scale. The defense is Schelling stickiness:
participants who have converged on a focal point typically hold
through attempts to move the focus, because moving costs
coordination value.

\paragraph{Dump attacks.}
An adversary with large \fig{} holdings can attempt to
destabilize focal-point status by dumping supply. In a pure
Schelling-point frame, this is \emph{less damaging} than in a
yield-anchored or governance-anchored frame: holders of \fig{}
are holders by alignment choice, not by financial hedging, and
the financial price of \fig{} is somewhat decoupled from its
focal-point status. A price crash that leaves participant
alignment unchanged does not necessarily destroy the focal
point; a coordinated loss of participant alignment does.

\paragraph{Acquiescence attacks.}
The quietest attack is no attack at all: nobody converges, and
\fig{} fails to become a focal point. This is possible and
unavoidable. The defense is limited --- we can support focal
emergence but cannot force it. The honest statement is that if
no cohort recognizes \fig{} as a focal point, the token has no
Schelling-point value and operates as an unvalued distribution
artifact. The supply-integrity and allocation properties
continue to hold in that case; only the signaling function
fails.

\paragraph{The griefing attack reconsidered.}
In earlier design iterations, a concern was raised about
\emph{griefing} --- an adversary who writes code to drain the
reward pool not for profit but for spite, given that no
governance or fee-extraction capture is available. Under Euler,
this attack had a concrete surface (manipulate the emission
function to empty the reward budget). Under the current
pure-airdrop design, the attack surface is narrower: the
airdrop contract's claim path requires merkle-proof inclusion
at a pre-committed address, and no drain is possible beyond
what the merkle root already authorizes. The
dump-at-unlock surface remains (addresses can claim at unlock
and sell), but that is indistinguishable from the normal
operation of the airdrop; it is not an exploit.

% ════════════════════════════════════════════════════════════════════
\section{Substrate-Neutrality and Distributional Silhouette}
\label{sec:silhouette}

The bonded settlement \emph{kernel} is a substrate on which
multiple organizational forms can be expressed. The present
paper describes \fig{}, a specific token ecosystem that ships
alongside the kernel. Those two statements are about different
objects and a careful reader should not confuse them.

\fig{}'s 10/30/60 allocation is \emph{one} instantiation of a
token economy over the substrate. The substrate is
substrate-neutral; \fig{} is not. If the distributional
silhouette of \fig{} --- early founder allocation, large
standing treasury, community recipients on a calendar schedule
--- rhymes with the silhouette of prior platform-capital
arrangements that the political-economy paper critiques, that
rhyme is a feature of this specific token's distribution, not
of the substrate, and not inherent to the Schelling-point
framing. Different focal-point tokens could be launched with
different silhouettes (flatter, community-only, non-founder),
and the Schelling-point design principles would apply to them
unchanged.

We flag this explicitly to disentangle three claims:

\begin{itemize}
  \item \textbf{Substrate-neutrality (true of the kernel)}: the
    bonded settlement primitive does not impose any particular
    distributional structure on tokens or organizations.
  \item \textbf{Focal-point integrity (true of \fig{})}: the
    specific supply mechanism and exclusion catalogue enforce
    the allocation and the focality-preservation properties at
    deploy.
  \item \textbf{Allocation legitimacy (out-of-scope)}: whether
    10/30/60 is the \emph{right} allocation, whether a 30\%
    liquid treasury is an appropriate ecosystem-coordination
    vehicle, and whether the implicit authority of the DAO over
    grant direction reproduces the soft-governance patterns
    platform-hegemony analysis critiques --- these are
    political-economy questions the present paper does not
    resolve.
\end{itemize}

% ════════════════════════════════════════════════════════════════════
\section{Verification Surface}
\label{sec:verification}

The \fig{} ecosystem is verified independently of the kernel.
We report rule and invariant counts as of the snapshot date;
numbers will grow as the surface expands.

\paragraph{TLA\textsuperscript{+} (\texttt{formal/FigToken.tla}).}
8 invariants verified across 160{,}844 distinct states in
approximately 9 seconds (TLC, breadth-first search, parameters
$\texttt{MAX\_SUPPLY}=5$, 3 minters, 6 recipients):
\texttt{Inv\_MaxSupply}, \texttt{Inv\_NonNegative},
\texttt{Inv\_DeployerCannotMintAfterRenounce},
\texttt{Inv\_NoMintToZero}, \texttt{Inv\_MinterCap},
\texttt{Inv\_BalancesSumToSupply},
\texttt{Inv\_CapBelowMaxSupply},
\texttt{Inv\_SupplyEqualsSumMinted}.

\paragraph{Halmos (retired).}
The original verification round included a 4-property
symbolic-execution harness
(\texttt{test/HalmosStagedMerkleAirdrop.t.sol}) covering
claim flag set on success, one-shot enforcement per (stage,
address), correct balance arithmetic, and merkle leaf format
$\texttt{keccak256(abi.encodePacked(account, amount))}$. The
harness was retired in 2026-05 alongside the
\texttt{StagedMerkleAirdrop} contract; the replacement
\texttt{RpgfMinter} is covered by Foundry tests and a
Rust\,$\leftrightarrow$\,Solidity conformance harness but does
not yet carry a Halmos harness of its own.

\paragraph{Certora.}
6 CVL rules in \texttt{certora/FigToken.spec}:
\texttt{totalSupplyWithinMaxSupply},
\texttt{totalRegisteredCapWithinMaxSupply},
\texttt{minterMintedWithinCap},
\texttt{totalRegisteredCapMonotonic},
\texttt{deployerMintRenouncedIsOneWayLatch},
\texttt{minterCapImmutable}.

A 3-rule spec for the staged airdrop contract
(\texttt{certora/StagedMerkleAirdrop.spec}: \texttt{claimedIsMonotonic},
\texttt{stageConfigImmutable}, \texttt{minterImmutable}) was part of
the original verification round but was retired in 2026-05 when
the contract was replaced by \texttt{RpgfMinter.sol}; the
replacement does not yet carry a Certora spec.

\paragraph{Foundry regression.}
\texttt{test/fig/FigToken.t.sol}, \texttt{test/fig/RpgfMinter.t.sol},
and \texttt{test/fig/RpgfMinterConformance.t.sol} cover the deploy
flow, every custom-error revert path, the
Rust\,$\leftrightarrow$\,Solidity conformance assertion on the
canonical merkle root, and the canonical 10/30/60 allocation realized
by \texttt{script/DeployMainnet.s.sol}. (The earlier
\texttt{test/fig/StagedMerkleAirdrop.t.sol} suite was retired
2026-05 alongside the contract.)

\paragraph{Headline supply guarantee.}
The verification surface yields a strong guarantee: from the
moment \texttt{deployerMintRenounced = true} is included in a
finalized transaction, the \texttt{totalSupply()} trajectory is
bounded above by $10^9$, no new minters can be registered, and
the only remaining mint authority is the airdrop contract whose
schedule is itself frozen by its three immutable stages. The
supply is deterministic conditional on the merkle leaves; the
merkle leaves are committed at deploy. Focal-point participants
assessing \fig{} can verify this directly from on-chain state.

% ════════════════════════════════════════════════════════════════════
\section{Regulatory Framing}
\label{sec:regulatory}

A token paper that takes itself seriously must address the
regulatory-status question, even if the answer it produces is
provisional and jurisdictionally bounded. We address the
question here in three registers --- US securities law,
EU MiCA, and a structural register prior to specific regimes.

\paragraph{The Howey four-prong test.}
Under \emph{SEC v. W.J. Howey Co.}, 328 U.S. 293 (1946) and its
extensive subsequent application to token offerings, an investment
contract exists where there is (i)~an investment of money,
(ii)~in a common enterprise, (iii)~with an expectation of profit,
(iv)~derived from the entrepreneurial or managerial efforts of
others. The SEC's 2019 \emph{Framework for Investment Contract
Analysis of Digital Assets} extended the test to digital-asset
arrangements specifically \citep{sec2019framework}, weighing
factors that bear on each prong. \fig{}'s position relative to
the four prongs is, on our reading, the following.

(i)~\emph{Investment of money.} A participant acquires \fig{}
through a merkle-claim airdrop or through secondary-market
purchase. The airdrop is not an investment of money in the
ordinary sense (no consideration is delivered to the issuer);
the secondary purchase is an investment of money in the same
sense any token acquisition is. Prong (i) is satisfied for
secondary acquisition and not for airdrop receipt. Most analyses
treat this asymmetry as not dispositive --- the secondary market
is the more material context --- so we proceed assuming
prong (i) is met for the relevant population.

(ii)~\emph{Common enterprise.} The doctrinal test for ``common
enterprise'' is jurisdictionally split (horizontal vs. vertical
commonality) but typically asks whether the investor's fortunes
are tied to the enterprise's fortunes through pooling of assets
or sharing of profits. \fig{} is not pooled, has no
profit-sharing arrangement, and is not tied to a treasury whose
performance distributes to holders. The DAO treasury holds 30\%
of supply (\Cref{sec:allocation}) and disburses grants under
its own governance, but treasury distributions are not pro-rata
to \fig{} holdings and do not constitute profit-sharing. Prong
(ii) is the weakest of the four for \fig{}; we read it as
likely \emph{not} satisfied under either commonality test for
\fig{} as currently constituted, with the caveat that the DAO's
disbursement policy is the relevant variable and a sufficiently
mechanistic disbursement policy could re-implicate the prong.

(iii)~\emph{Expectation of profit.} Holders may reasonably
expect \fig{} to appreciate if the focal convergence
strengthens. This is an expectation of profit in the ordinary
sense, and prong (iii) is met.

(iv)~\emph{Derived from the efforts of others.} This is the
prong on which most token cases turn under the SEC framework.
The framework weighs whether an active participant (developer,
foundation, promoter) is undertaking essential managerial efforts
on which the holder's profit expectation depends. \fig{}'s
deploy script registers the deployer as a one-shot minter,
mints the founder and DAO allocations, and renounces ---
permanently, by code-enforced one-way latch
(\Cref{sec:supply}). After deploy, the founder's effort is not
managerial in any sense the framework recognizes: the founder
holds liquid \fig{} as any holder does, has no protocol-level
authority, cannot mint, cannot upgrade, cannot freeze, and has
no fiduciary relationship to other holders. The DAO has
treasury authority over its own allocation but not protocol
authority over \fig{}. The development of the underlying
\figaro{} protocol is a separate question, governed by a
separate and ongoing analysis; \fig{} is not coupled to that
development through any protocol-level mechanism. Prong (iv)
is, on our reading, partially met (insofar as ecosystem
development is happening and holder profit expectations may
depend on it) and partially not met (insofar as the
relationship between holder profit and developer effort is
reflexive, not contractual or fiduciary). The prong is the
hardest to assess because the SEC framework's factors were
written with team-managed token projects in mind, and \fig{}
is structurally a less-managed instrument than the framework
contemplates.

The honest summary is that \fig{}'s position relative to Howey
is uncertain and varies with the development trajectory of the
broader \figaro{} ecosystem. We do not claim \fig{} is not a
security; we claim that the design is engineered to minimize
the prongs' applicability, and that the resulting position is
defensible but not assured. This is the most we believe a
cryptoeconomic paper can responsibly say. Specific
jurisdictional advice belongs in a regulatory submission, not
in this venue.

\paragraph{EU MiCA.}
The Markets in Crypto-Assets Regulation (Regulation (EU)
2023/1114) classifies crypto-assets into asset-referenced
tokens, e-money tokens, and ``other'' crypto-assets, with
distinct disclosure and authorization regimes. \fig{} is not
asset-referenced (no reserve), not e-money (no nominal-value
stability target), and would fall in the residual ``other''
category, subject to the white-paper publication requirement
of MiCA Title II. The MiCA regime is more permissive of
focal-point-style instruments than Howey is, conditional on
disclosure compliance, and we treat MiCA as a tractable
disclosure obligation rather than as an existential
classification question.

\paragraph{Structural register.}
A reader may ask why a paper that has spent itself arguing for
\fig{} as a focal-point coordination instrument should
nonetheless engage securities law at all. The answer is that
``what \fig{} is for'' (focal-point coordination) and ``what
\fig{} is in regulatory-classification space'' (a question
about the prongs) are different questions, and a paper that
addresses the first without acknowledging the second is doing
its readers a disservice. The DAO-allocation contradiction
flagged in \Cref{sec:allocation} is, in part, a regulatory
contradiction: a 30\% treasury with grant-disbursement
authority is exactly the structural feature most likely to
re-implicate prong (iv) of Howey if the disbursement policy
becomes mechanically tied to holding behavior. Keeping the
disbursement policy at human discretion, separated from
holding by the two layers of friction we identified there
(off-chain discretion, revocable grants), is also a regulatory
discipline as well as a focality-protection discipline. The
two disciplines align, and we have written the design to
satisfy both, but the alignment is not coincidental and should
not be relied on as a substitute for legal advice in any
specific jurisdiction.

% ════════════════════════════════════════════════════════════════════
\section{Conclusion}
\label{sec:conclusion}

\fig{} is a small design. Most of its positive content is
procedural (supply cap, renouncement latch, staged merkle
airdrop) and most of its substantive content is negative (no
governance, no yield, no fee path, no founder vesting, no
settlement-anchored emission, no coordinator-pattern role).
Under the framing this paper develops, the negatives are not
restraint without purpose --- they are the specific design
discipline required to support focal-point emergence. Every
exclusion removes a class of holder whose presence would dilute
the signaling value of \fig{} for participants who hold for
alignment reasons.

\paragraph{The lineage, made explicit.}
\citet{nakamoto2008bitcoin} took money out of institutions and
made it a protocol. Bitcoin is the Schelling point at the money
layer, underpinned by the cost of producing the asset: a bearer
asset around which participants converge because its design
makes convergence credible (fixed supply, issuer-less,
peer-to-peer). \citet{buterin2014ethereum} made contracts
composable and put \textsc{Eth} at the focal point of
programmable settlement, underpinned by transaction demand: gas
is the utility floor, Schelling recognition the coordination
premium. The present work extends the pattern to institutions:
bonded commitments are composable into arbitrarily large
assemblies (mechanism and implementation papers), and
\fig{} is the focal point for aligned participation in that
composition, underpinned by \emph{trade} --- the value-added
produced when parties exchange goods, services, or coordination
through the bonded primitive. The pattern is the same at each
layer: a native token that is not a utility, not a governance
claim, and not a cash-flow instrument, but a coordination anchor
whose value rests on a distinct real-economic quantity. The
distinction between the three is which quantity: factor cost for
Bitcoin, transaction utility for Ethereum, and the value-added
of trade for \fig{}.

\paragraph{Open questions.}
Three questions remain outside this paper's scope but worth
flagging. First, the airdrop cohort selection (which addresses
appear in which stage's merkle root) is policy, not mechanism,
and the policy framework should be developed and audited as a
separate artifact before each stage's root is committed at
deploy. Second, the relationship between \fig{} and the various
treasury communities likely to form around it (cf.\ the
``capital firm'' discussion in the institutional-economics paper) is
genuinely emergent and not centrally designable; the present
paper does not attempt to predict or control that emergence.
Third, the status of the specific percentages in the allocation
is as addressed in \Cref{sec:allocation}: judgment, not proof;
the structural choices are the defensible content.

\paragraph{What the design can and cannot do.}
Schelling-point design is a matter of supporting emergence, not
causing it. The design rules of \Cref{sec:design-rules} and the
exclusions of \Cref{sec:negatives} make focal-point convergence
on \fig{} viable: a participant who wants to signal alignment
can do so by acquiring and holding \fig{}, and the design
ensures the signal is not muddied by non-alignment holders.
Whether focal-point convergence \emph{happens} is a social
process we do not control. We document the design so that those
who choose to converge can do so legibly, and so that those who
do not can see what they are declining.

\section*{Acknowledgments}
The decision to frame \fig{} as a Schelling-point token,
rather than as a decoupled-emission token whose cryptoeconomic
content is largely negative, emerged in discussions with
collaborators who pressed the question of what \fig{}
\emph{is} rather than what it \emph{is not}. Their skepticism
clarified the argument. The verification surface owes a
continuing debt to the maintainers of TLA\textsuperscript{+},
Halmos, and Certora.


% ════════════════════════════════════════════════════════════════════
% References
% ════════════════════════════════════════════════════════════════════

\begin{thebibliography}{99}

\bibitem[Buterin(2014)]{buterin2014ethereum}
Buterin, V.
\newblock Ethereum: A Next-Generation Smart Contract and
  Decentralized Application Platform.
\newblock White paper, 2014.
\newblock \url{https://ethereum.org/en/whitepaper/}

\bibitem[Nakamoto(2008)]{nakamoto2008bitcoin}
Nakamoto, S.
\newblock Bitcoin: A Peer-to-Peer Electronic Cash System.
\newblock White paper, 2008.
\newblock \url{https://bitcoin.org/bitcoin.pdf}

\bibitem[Sugden(1995)]{sugden1995theory}
Sugden, R.
\newblock A Theory of Focal Points.
\newblock \emph{Economic Journal}, 105(430):533--550, 1995.

\bibitem[Bacharach(2006)]{bacharach2006beyond}
Bacharach, M.
\newblock \emph{Beyond Individual Choice: Teams and Frames in
  Game Theory}, edited by N.~Gold and R.~Sugden.
\newblock Princeton University Press, Princeton, NJ, 2006.

\bibitem[Mehta et al.(1994)]{mehta_starmer_sugden_1994}
Mehta, J., Starmer, C., and Sugden, R.
\newblock The Nature of Salience: An Experimental Investigation
  of Pure Coordination Games.
\newblock \emph{American Economic Review}, 84(3):658--673, 1994.

\bibitem[Markey-Towler(2018)]{markey-towler2018schelling}
Markey-Towler, B.
\newblock Anarchy, Blockchain and Utopia: A Theory of Political-Economic
  Systems and the Coordination of Blockchain Networks.
\newblock \emph{Journal of the British Blockchain Association},
  1(1):1--15, 2018.

\bibitem[Daian et al.(2017)]{daian_etal_2017}
Daian, P., Pass, R., and Shi, E.
\newblock On the Instability of Bitcoin Without the Block Reward.
\newblock In \emph{Proceedings of the 2016 ACM SIGSAC Conference on
  Computer and Communications Security}, 154--167, 2016.

\bibitem[Buterin(2020)]{buterin2020credible}
Buterin, V.
\newblock Credible Neutrality As A Guiding Principle.
\newblock Personal blog, 2020.
\newblock \url{https://nakamoto.com/credible-neutrality/}

\bibitem[Hayek(1945)]{hayek1945knowledge}
Hayek, F.~A.
\newblock The Use of Knowledge in Society.
\newblock \emph{American Economic Review}, 35(4):519--530, 1945.

\bibitem[Selgin(1988)]{selgin1988theory}
Selgin, G.~A.
\newblock \emph{The Theory of Free Banking: Money Supply Under
  Competitive Note Issue}.
\newblock Rowman \& Littlefield, Totowa, NJ, 1988.

\bibitem[SEC(2019)]{sec2019framework}
US Securities and Exchange Commission.
\newblock \emph{Framework for ``Investment Contract'' Analysis of
  Digital Assets}.
\newblock SEC Strategic Hub for Innovation and Financial Technology,
  April 3, 2019.

\bibitem[Schelling(1960)]{schelling1960strategy}
Schelling, T.~C.
\newblock \emph{The Strategy of Conflict}.
\newblock Harvard University Press, Cambridge, MA, 1960.

\end{thebibliography}

\end{document}
